\documentclass[12pt]{article}
\usepackage[utf8]{inputenc}
\usepackage[T1]{fontenc}
\usepackage{hyperref}
\title{Muruwari Language Resources for K-12 Curriculum Development}
\author{Compiled by Claude Sonnet 5.5 (claude-sonnet-5-5) for Andr\'as Kornai}
\date{}
\begin{document}
\maketitle

\section{Introduction}
Muruwari (Murawari, Muruwarri; ISO 639-3 zmu, Glottocode muru1266) is a dormant Pama-Nyungan language of northwestern New South Wales and the Goodooga, Brewarrina and Bourke area. Most of the surviving material comes from recordings made by Jimmie Barker of Brewarrina and from Lynette Oates's work with the last speakers in the 1980s and 1990s. Wikipedia reports about 12 people with some revival ability (2016).

\section{Primary Educational Resources}
No school primers or textbooks were found. Revival work is community-based, and the sources below describe it.
\begin{itemize}
\item The Western Herald on NSW funding for Aboriginal language revitalisation, including the Brewarrina Local Aboriginal Lands Council's First Linguists workshops covering Muruwari and other languages: \url{https://www.thewesternherald.com.au/single-post/4-million-to-revitalise-aboriginal-languages}
\end{itemize}

\section{Linguistic and Reference Materials}
\begin{itemize}
\item AIATSIS Austlang entry D32: \url{https://collection.aiatsis.gov.au/austlang/language/D32}
\item Wikipedia on the Muruwari language: \url{https://en.wikipedia.org/wiki/Muruwari_language}
\item Wikipedia on the Muruwari people: \url{https://en.wikipedia.org/wiki/Muruwari}
\item The Language Closet overview: \url{https://thelanguagecloset.com/2022/01/22/languages-of-australia-muruwari/}
\end{itemize}

\section{Religious and Community Resources}
The State Library of Queensland says that community revival is supported by several organisations, including Wonganurra Aboriginal Corporation and Pulkurru Aboriginal Corporation: \url{https://www.slq.qld.gov.au/blog/2019-international-year-indigenous-languages-word-week-week-forty-three}

\section{Digital and Audio Resources}
A federal grant funded documentation of the remaining Jimmie Barker records, intended to support an audio-based online dictionary and a podcast series. I did not confirm whether these have been published. See the Western Herald article above, and the Rosetta Panglossia entry: \url{https://rosettapanglossia.longnow.org/index.php?title=Zmu}

\section{Curriculum Development Context}
The language is dormant and not known to be taught in schools. A heritage and revitalisation curriculum would rely on the Barker and Oates documentation, and would be developed with community groups around Brewarrina, Bourke and Goodooga.
\end{document}